% ============================================================
% Lists.tex — Abbreviations and Symbols
%
% HOW TO ADD AN ABBREVIATION:
%   \newacronym{KEY}{SHORT}{Full Form}
%   Example: \newacronym{api}{API}{Application Programming Interface}
%
%   In your chapter text, use \gls{api} — on first use it
%   expands to "Application Programming Interface (API)".
%   On subsequent uses it shows just "API".
%
% HOW TO ADD A SYMBOL:
%   \newglossaryentry{KEY}{type=symbols, name={SYMBOL},
%     description={what it means}, sort={sort key}}
%
% If you are not using abbreviations or symbols, leave this
% file empty or delete the \printglossary lines in Thesis.tex.
% ============================================================


% ============================================================
% ABBREVIATIONS
% Replace the examples below with your own.
% ============================================================

\newacronym{ai}{AI}{Artificial Intelligence}
\newacronym{ml}{ML}{Machine Learning}
\newacronym{dl}{DL}{Deep Learning}
\newacronym{api}{API}{Application Programming Interface}
\newacronym{ui}{UI}{User Interface}
\newacronym{ux}{UX}{User Experience}
\newacronym{os}{OS}{Operating System}
\newacronym{sql}{SQL}{Structured Query Language}
\newacronym{rest}{REST}{Representational State Transfer}
\newacronym{gpu}{GPU}{Graphics Processing Unit}
\newacronym{cpu}{CPU}{Central Processing Unit}
\newacronym{ide}{IDE}{Integrated Development Environment}

% Add your own:
% \newacronym{KEY}{SHORT}{Full Form}


% ============================================================
% SYMBOLS (uncomment and edit if needed)
% ============================================================

% \newglossaryentry{n}{
%     type=symbols,
%     name={\ensuremath{n}},
%     description={Number of samples or observations},
%     sort={n}
% }

% \newglossaryentry{alpha}{
%     type=symbols,
%     name={\ensuremath{\alpha}},
%     description={Learning rate or significance threshold},
%     sort={alpha}
% }

% \newglossaryentry{bigO}{
%     type=symbols,
%     name={\ensuremath{\mathcal{O}}},
%     description={Big-O notation for algorithmic complexity},
%     sort={O}
% }
